\section{A date range can leave several information views possible}
\label{sim:exact-temporal-views}

A filing marked as published on 2 January does not establish whether a manager had read it at the beginning of that day. Nor does a correction received within a narrow interval around a decision establish whether the original or corrected interpretation was available for that decision. These are questions about the order of events. Assigning midnight or the interval's midpoint would settle them by assumption.

The selected successor retains that uncertainty explicitly. An exact timestamp is a point; a calendar day in UTC is a half-open interval from its midnight to the next midnight. A supplied time range has the same half-open convention. An unknown access time remains unknown. Exact rational seconds preserve differences smaller than a microsecond and equivalent fixed-offset spellings without passing through a floating-point or microsecond clock.

If the access time lies in $[a,b)$, a decision at $b$ is certainly after access. A decision at $a$ can coincide with access or precede it: the lower endpoint belongs to the interval. The distinction disappears if an implementation replaces the interval with a single representative time. It also disappears if a strict upper bound is inadvertently changed to a closed one.

\subsection{Availability depends on the route and on joint chronology}

A public-document retrieval and a private delivery impose different conditions. In the first route the source must have been released publicly before the claimed access. In the second, an actor-addressed delivery can establish access before any public release. The private route neither invents a public timestamp nor makes the document available to other actors. Both routes still require evidence about that actor's access; publication alone cannot supply it.

The earlier whole-second reference did not distinguish these routes. One of its retained cases places actor access before a later public release. Under the new, explicit public-retrieval premise, that chronology is inconsistent. It can be meaningful under a private-delivery premise. Preserving the difference is necessary: changing the access mechanism changes the question rather than improving timestamp precision alone.

Clocks are also dependent. Authorship precedes ingestion, ingestion precedes review, and authorship precedes access. Public retrieval additionally follows publication. The selected revision profile requires nondecreasing authorship, ingestion, review and same-actor access along a replacement chain, while permitting retrospective effective dates. These are declared premises of this bounded source model. A document practice that violates one of them requires another profile; a corrected economic effective date is not forced forward merely to preserve a convenient time ordering.

An unknown clock need not make every conclusion unknown. If a private access time is supplied before the decision and the document was necessarily authored before that access, the chronology can establish the relevant authorship inequality without identifying its timestamp. The result must retain the unknown-time marker. Logical implication from other premises does not turn an unobserved time into an observed fact.

\subsection{Partition clock possibilities without enumerating timestamps}

Let $x_e$ be the time of event $e$, measured in exact UTC seconds on the retained calendar axis. Point and interval observations, chronology and cutoff conditions become difference inequalities. A non-strict edge represents
\begin{equation}
 x_v-x_u\leq c,
 \label{eq:temporal-difference-edge}
\end{equation}
with a strict variant when the right endpoint is excluded. A fixed origin converts an absolute bound into the same form. The source identities, actors, access routes and original precision remain attached to these mathematical restrictions; the inequalities alone do not establish the source facts.

For example, publication before access contributes $x_{\mathrm{publication}}\leq x_{\mathrm{access}}$, while access by the decision contributes $x_{\mathrm{access}}\leq t_{\mathrm{decision}}$. Keeping both inequalities in one system prevents the calculation from accepting a combination that separate per-clock classifications would mistakenly permit.

For each revision, eligibility is a conjunction of its required actor, system and effective-time guards, $g_1,\ldots,g_m$. The complementary cases can be written as a disjoint first-failure partition:
\begin{equation}
 \neg(g_1\land\cdots\land g_m)
 =\neg g_1\;\lor\;(g_1\land\neg g_2)\;\lor\;\cdots\;\lor\;
 (g_1\land\cdots\land g_{m-1}\land\neg g_m).
 \label{eq:temporal-first-failure}
\end{equation}
Each branch preserves the earlier satisfied guards. Repeating the construction for the relevant revisions covers every clock assignment admitted by the original ranges and chronology. Eligibility and the selected revision heads are constant within each surviving branch. The procedure therefore represents families of timestamps symbolically; it does not need an explicit tree point for every possible instant.

An empty branch requires a checked contradiction. Summing original difference inequalities around a cycle cancels all clock variables. A negative right-hand sum is impossible; a zero sum is also impossible if at least one inequality in the sum is strict. A nonempty branch instead carries exact rational clock values satisfying every original inequality. The checker verifies those values directly and reconstructs every complementary branch before accepting the partition.

The witness establishes nonemptiness, while the guards and complete partition establish the scope of the conclusion. A few feasible timestamp examples would establish neither exhaustive coverage nor the absence of an omitted information view. Conversely, reaching the computation budget leaves an unfinished temporal family. The selected reference then reports incompleteness and does not present the already examined views as the full set.

\subsection{Different timing views can support different policies}

Return to receipt 7 and the committed facility. Suppose a correction changes the available observation scheme from coarse to refined, but the actor's access interval straddles the first decision. There are two remaining timing views. In the view where the correction has not yet arrived, the coarse signal leaves incompatible first instructions in the same information cell and no robust policy is feasible. In the view where it has arrived, the refined signal supports the four previously derived policies. The complete financial row population is retained in both calculations.

The timing evidence supplies no probability for either view. Selecting the four-policy answer would discard the other possibility. Reporting the two answers separately is therefore a substantive conclusion. Constructing one decision rule that is admissible when even the observation mechanism is uncertain requires a common specification of what the actor sees and how actions are represented across those mechanisms; a naive intersection of differently labelled policy tables is insufficient.

The exact comparison verifies interval boundaries, equivalent offsets, unknown and untrusted clocks, public and private access, correlated correction times, conflicting revisions and retractions. The saved partitions retain unresolved source outcomes as well as selected ones. Their scope is the supplied chronology and access model, not an empirical statement about how a real company releases information. General revision practices, authenticated evidence and durable recovery remain separate questions; Section~\ref{sim:temporal-family-policies} develops complete-family reuse and common decisions across possible observation mechanisms.\footnote{The derivation, original-edge checker, complete saved temporal partitions, independent finite order comparisons and full policy sets are retained in \path{data/accounting-master/source-time-intervals-20261027/}. The reference preserves its failed namespace diagnostic, the explicit public-route difference from the whole-second profile, and the six-revision and computation-budget limits. No clock probabilities or actual actor-access admission are introduced.}
